% ============================================================
%  Příprava na maturitu – Mechatronika
%  Kompiluj: pdflatex maturita_mechatronika.tex
%  Kódování: UTF-8
% ============================================================

\documentclass[a4paper,10pt]{article}

% --- Balíčky ---
\usepackage[T1]{fontenc}
\usepackage[utf8]{inputenc}
\usepackage[czech]{babel}
\usepackage[left=2.4cm, right=1.8cm, top=1.8cm, bottom=1.8cm]{geometry}
\usepackage{lmodern}
\usepackage{microtype}
\usepackage{parskip}
\usepackage{titlesec}
\usepackage{enumitem}
\usepackage{xcolor}
\usepackage{hyperref}
\usepackage{tabularx}
\usepackage{booktabs}
\usepackage{array}
\usepackage{tikz}
\usepackage{eso-pic}
\usepackage{mdframed}
\usepackage{amsmath}
\usepackage{amssymb}
\usepackage{pgfplots}
\pgfplotsset{compat=1.18}
\usetikzlibrary{arrows.meta, positioning, shapes.geometric, calc}

% --- Barvy ---
\definecolor{accent}{HTML}{03254E}
\definecolor{stripeblue}{HTML}{568EA3}
\definecolor{medgray}{HTML}{888888}
\definecolor{lightblue}{HTML}{EAF2F8}
\definecolor{lightyellow}{HTML}{FDFBE4}
\definecolor{lightgreen}{HTML}{EAF7EA}

% --- Modrý pruh vlevo ---
\AddToShipoutPictureBG{%
  \begin{tikzpicture}[remember picture, overlay]
    \fill[stripeblue]
      (current page.north west)
      rectangle
      ([xshift=10mm] current page.south west);
    \fill[accent!60]
      ([xshift=10mm] current page.north west)
      rectangle
      ([xshift=12mm] current page.south west);
  \end{tikzpicture}%
}

\hypersetup{colorlinks=true, urlcolor=accent, linkcolor=accent}

\titleformat{\section}
  {\large\bfseries\color{accent}}
  {\thesection.}
  {0.5em}
  {}
  [\color{accent}\titlerule]

\titlespacing{\section}{0pt}{16pt}{6pt}

\newmdenv[
  backgroundcolor=lightblue,
  linecolor=stripeblue,
  linewidth=1pt,
  roundcorner=4pt,
  innerleftmargin=8pt,
  innerrightmargin=8pt,
  innertopmargin=6pt,
  innerbottommargin=6pt,
  skipabove=4pt,
  skipbelow=4pt
]{pojmybox}

\newmdenv[
  backgroundcolor=lightyellow,
  linecolor=accent!40,
  linewidth=1pt,
  roundcorner=4pt,
  innerleftmargin=8pt,
  innerrightmargin=8pt,
  innertopmargin=6pt,
  innerbottommargin=6pt,
  skipabove=4pt,
  skipbelow=8pt
]{vysvetlenibox}

\newmdenv[
  backgroundcolor=lightgreen,
  linecolor=accent!30,
  linewidth=1pt,
  roundcorner=4pt,
  innerleftmargin=8pt,
  innerrightmargin=8pt,
  innertopmargin=6pt,
  innerbottommargin=6pt,
  skipabove=4pt,
  skipbelow=8pt
]{vzorcebox}

\pagestyle{plain}

% ============================================================
\begin{document}

% --- Záhlaví ---
\begin{minipage}[t]{0.75\textwidth}
    {\Huge\bfseries\color{accent} Příprava na maturitu}\\[4pt]
    {\large\color{stripeblue} Mechatronika}\\[2pt]
    \textcolor{medgray}{\small Suchánek Lukáš \quad SPŠ Prosek \quad 2026}
\end{minipage}
\hfill
\begin{minipage}[t]{0.22\textwidth}
    \raggedleft
    \begin{tikzpicture}
        \draw[color=stripeblue, rounded corners=4pt, line width=1pt]
            (0,0) rectangle (3,1.2)
            node[midway, align=center, color=accent, font=\small\bfseries]
            {Otázek: 22};
    \end{tikzpicture}
\end{minipage}

\vspace{6pt}
\noindent\color{accent}\rule{\linewidth}{1.5pt}
\vspace{2pt}

% ============================================================
\section{CAx systémy a Rapid prototyping}

\begin{pojmybox}
\textbf{Klíčové pojmy:}
CAD, CAM, CAE, FEM, simulace, Fusion~360, SolidWorks, CATIA,
Rapid prototyping, FDM, SLA, SLS, DMLS, postprocesor,
STL formát, topologická optimalizace
\end{pojmybox}

\begin{vysvetlenibox}
\textbf{CAx -- přehled nástrojů}

\begin{center}
\begin{tabularx}{0.95\linewidth}{@{} l l X @{}}
    \toprule
    \textbf{Zkratka} & \textbf{Název} & \textbf{Funkce} \\
    \midrule
    CAD & Computer Aided Design      & 3D parametrické modelování geometrie \\
    CAM & Computer Aided Manufacturing & Generování obráběcích drah z modelu \\
    CAE & Computer Aided Engineering & Simulace a analýza (FEM, CFD, MBS) \\
    CAPP& Aided Process Planning     & Plánování výrobního procesu \\
    PDM & Product Data Management    & Správa dat a verzí modelu \\
    \bottomrule
\end{tabularx}
\end{center}

\textbf{Metoda konečných prvků (FEM)}\\
Soustava se rozdělí na malé prvky (mesh). Pro každý prvek se zapíše
rovnice tuhosti. Soustava rovnic se řeší numericky.

\begin{vzorcebox}
Maticová forma: $[K]\{u\} = \{F\}$\\
kde $[K]$ = globální matice tuhosti, $\{u\}$ = posuvy uzlů,
$\{F\}$ = vektor sil.\\
Napětí: $\sigma = E \cdot \varepsilon = E \cdot \frac{\Delta l}{l_0}$
\end{vzorcebox}

\textbf{Rapid prototyping -- pracovní postup}

\begin{center}
\begin{tikzpicture}[scale=0.85,
  step/.style={rectangle, draw=accent, fill=lightblue, rounded corners=3pt,
               minimum width=2.3cm, minimum height=0.6cm,
               text centered, font=\small}]
  \node[inner sep=2pt,rectangle,draw,fill=blue!20] (cad)  at (0,0)    {3D model (CAD)};
  \node[inner sep=2pt,rectangle,draw,fill=blue!20] (stl)  at (3,0)    {Export STL};
  \node[inner sep=2pt,rectangle,draw,fill=blue!20] (slic) at (6,0)    {Slicování};
  \node[inner sep=2pt,rectangle,draw,fill=blue!20] (tisk) at (6,-1.2) {Tisk / výroba};
  \node[inner sep=2pt,rectangle,draw,fill=blue!20] (post) at (3,-1.2) {Post-processing};
  \node[inner sep=2pt,rectangle,draw,fill=blue!20] (test) at (0,-1.2) {Testování};
  \draw[>-,accent] (cad)--(stl)--(slic)--(tisk)--(post)--(test);
\end{tikzpicture}
\end{center}

\textbf{Topologická optimalizace}\\
Algoritmus odstraní materiál z míst s nízkým napětím -- výsledek
je organická struktura s maximálním poměrem tuhosti/hmotnosti.
Typicky se kombinuje s DMLS tiskem (nelze vyrobit klasickým obráběním).

\textbf{Srovnání RP technologií -- použití v praxi}

\begin{center}
\begin{tabularx}{0.95\linewidth}{@{} l X X @{}}
    \toprule
    \textbf{Metoda} & \textbf{Výhody} & \textbf{Nevýhody} \\
    \midrule
    FDM  & Levná, dostupná, velký sortiment materiálů
         & Viditelné vrstvy, nižší přesnost \\
    SLA  & Hladký povrch, vysoká přesnost
         & Drahá pryskyřice, křehké \\
    SLS  & Bez podpor, funkční díly z nylonu
         & Drahý stroj, prašné \\
    DMLS & Kovové díly, složité tvary, vysoká pevnost
         & Velmi drahý stroj (\textgreater{}500k EUR) \\
    \bottomrule
\end{tabularx}
\end{center}
\end{vysvetlenibox}

% ============================================================
\vfill
\noindent\textcolor{medgray}{\small
Souhlasím se zpracováním osobních údajů pro účely náboru v souladu s GDPR.}

\end{document}
```

% ============================================================
\end{document}
